% =============================================================================
% CHAPTER 10: NARRATIVE AI STABILITY
% IP Traceability: IP-009 (narrariveai (1).pdf)
% =============================================================================

\chapter{Narrative AI Stability}
\label{chap:narrative-stability}

\section{IP Document Summary}

\begin{tracerecord}
\textbf{IP Source ID} & \ipid{IP-009} \\
\textbf{Document} & narrariveai (1).pdf \\
\textbf{Author} & Craig Turrell, Head of AI \& Design \\
\textbf{Date} & January 2026 \\
\textbf{Pages} & 18 \\
\textbf{Abstract} & Lyapunov stability metrics and narrative graph neural networks for AI governance \\
\end{tracerecord}

\subsection{Document Abstract}

\begin{ipsourcebox}{IP-009, Abstract}
\textit{``This paper introduces a formal stability framework for AI narrative generation systems. We apply Lyapunov stability theory to define convergence guarantees for narrative coherence, ensuring that AI-generated narratives remain bounded and predictable. The framework includes real-time stability dashboards, traffic-light status indicators, and integration with enterprise governance systems.''}
\end{ipsourcebox}

\section{Key Concepts Identified}

\begin{table}[H]
\centering
\caption{IP-009 Key Concepts}
\begin{tabularx}{\textwidth}{@{}p{3.5cm}Xl@{}}
\toprule
\textbf{Concept} & \textbf{Description} & \textbf{Section} \\
\midrule
Lyapunov Stability Metrics & Mathematical stability measurements & \ref{sec:ip009-lyapunov} \\
Stability Dashboard & Real-time monitoring interface & \ref{sec:ip009-dashboard} \\
Traffic-light Status & Intuitive status indicators & \ref{sec:ip009-traffic} \\
Governance Integration & TOON and metadata compliance & \ref{sec:ip009-governance} \\
\bottomrule
\end{tabularx}
\end{table}

\section{Concept: Lyapunov Stability Metrics}
\label{sec:ip009-lyapunov}

\begin{ipsourcebox}{IP-009, Stability Framework Section}
\textit{``We define the Maximum Lyapunov Exponent (MLE) for narrative trajectories:
$$\lambda = \lim_{t \to \infty} \frac{1}{t} \ln \frac{||\delta(t)||}{||\delta(0)||}$$
where $\delta(t)$ is the deviation from the expected narrative trajectory. Systems with $\lambda < 0$ are stable; $\lambda > 0$ indicates instability.''}
\end{ipsourcebox}

\begin{tracerecord}
\textbf{IP Source ID} & IP-009 \\
\textbf{IP Location} & Stability Framework Section \\
\textbf{Concept} & Quantitative stability measurement \\
\midrule
\textbf{Implementation Path} & Conceptual influence \\
\textbf{Adaptation Type} & Design principle (deviation tracking) \\
\end{tracerecord}

\begin{adaptbox}
The IP paper defines stability through Lyapunov exponents measuring deviation growth. While SAP-OSS does not compute MLE directly, the concept of tracking deviation from expected trajectories is implemented in anomaly detection: Z-scores and MAD scores measure how far values deviate from expected baselines. Large deviations (high scores) indicate ``instability'' in the data.
\end{adaptbox}

\section{Concept: Stability Dashboard}
\label{sec:ip009-dashboard}

\begin{ipsourcebox}{IP-009, Monitoring Section}
\textit{``The stability dashboard provides real-time visibility into AI system behavior:
\begin{itemize}
\item Current stability score (0--100)
\item Trend indicators (improving/stable/degrading)
\item Threshold alerts (warning/critical zones)
\item Historical stability charts
\end{itemize}
This enables proactive intervention before instability becomes critical.''}
\end{ipsourcebox}

\begin{tracerecord}
\textbf{IP Source ID} & IP-009 \\
\textbf{IP Location} & Monitoring Section \\
\textbf{Concept} & Real-time stability monitoring \\
\midrule
\textbf{Implementation Path} & \filepath{src/intelligence/analytics/} \\
\textbf{Adaptation Type} & Structural (monitoring data structures) \\
\end{tracerecord}

\begin{adaptbox}
The IP paper describes a real-time stability dashboard. The SAP-OSS implementation provides the data required for such monitoring through the result dataclasses: \funcname{BatchAnomalyResult} includes statistics (for current scores), method selection (for trend analysis), and threshold comparisons (for alerts). These structures support dashboard visualization.
\end{adaptbox}

\section{Concept: Traffic-light Status}
\label{sec:ip009-traffic}

\begin{ipsourcebox}{IP-009, Status Indicators Section}
\textit{``Traffic-light status indicators provide intuitive stability communication:
\begin{itemize}
\item \textbf{Green}: Stable operation ($\lambda < 0$, within normal bounds)
\item \textbf{Amber}: Warning ($\lambda \approx 0$, approaching thresholds)
\item \textbf{Red}: Critical ($\lambda > 0$, unstable, intervention required)
\end{itemize}
These map continuous metrics to discrete, actionable states.''}
\end{ipsourcebox}

\begin{tracerecord}
\textbf{IP Source ID} & IP-009 \\
\textbf{IP Location} & Status Indicators Section \\
\textbf{Concept} & Discrete status classification \\
\midrule
\textbf{Implementation Path} & \filepath{src/intelligence/analytics/anomaly\_detector.py} \\
\textbf{Adaptation Type} & Structural (anomaly\_type classification) \\
\end{tracerecord}

\begin{implbox}{Traffic-light Pattern in anomaly\_type}
\begin{lstlisting}[language=Python]
# Traffic-light status mapping in VarianceAnalysisResult
#
# anomaly_type provides discrete classification:
# - "none" -> Green (normal operation)
# - "statistical" -> Amber (statistical deviation detected)
# - "spike"/"drop" -> Amber/Red (significant directional change)
# - "material" -> Red (material threshold exceeded)
# - "new_account" -> Amber (requires attention)
#
# This maps continuous Z-scores to discrete, actionable states
# as specified in IP-009 traffic-light concept
\end{lstlisting}
\end{implbox}

\begin{adaptbox}
The IP paper defines traffic-light indicators mapping continuous metrics to discrete states. The SAP-OSS \texttt{anomaly\_type} field implements this concept: ``none'' indicates green (normal), ``statistical'' indicates amber (attention needed), and ``material''/``spike''/``drop'' indicate conditions requiring action. This provides the intuitive status communication described in the IP.
\end{adaptbox}

\section{Concept: Governance Integration}
\label{sec:ip009-governance}

\begin{ipsourcebox}{IP-009, TOON Integration Section}
\textit{``Stability metrics integrate with enterprise governance through Token-Oriented Object Notation (TOON):
\begin{itemize}
\item Stability scores included in TOON metadata
\item Audit trails link stability events to governance actions
\item Compliance reports include stability history
\end{itemize}
This ensures AI stability is part of overall governance posture.''}
\end{ipsourcebox}

\begin{tracerecord}
\textbf{IP Source ID} & IP-009 \\
\textbf{IP Location} & TOON Integration Section \\
\textbf{Concept} & Metadata-driven governance integration \\
\midrule
\textbf{Implementation Path} & \filepath{src/intelligence/ai-core-pal/governance/} \\
\textbf{Adaptation Type} & Conceptual (audit integration) \\
\end{tracerecord}

\begin{adaptbox}
The IP paper describes TOON metadata for governance integration. The SAP-OSS implementation captures this concept through the \texttt{details} dictionary in \funcname{VarianceAnalysisResult} and the structured audit records from \funcname{AuditLogger}. These provide the metadata required for compliance reports and governance dashboards, enabling stability tracking as part of overall AI governance.
\end{adaptbox}

\section{Implementation Summary}

\begin{table}[H]
\centering
\caption{IP-009 Concept Implementation Summary}
\begin{tabularx}{\textwidth}{@{}p{3cm}p{3.5cm}Xc@{}}
\toprule
\textbf{Concept} & \textbf{IP Reference} & \textbf{Implementation} & \textbf{Status} \\
\midrule
Lyapunov Metrics & Stability Section & Deviation tracking (Z-scores) & Yes \\
Stability Dashboard & Monitoring Section & Result dataclasses & Yes \\
Traffic-light Status & Status Section & \texttt{anomaly\_type} classification & Yes \\
Governance Integration & TOON Section & Audit logger + details dict & Yes \\
\bottomrule
\end{tabularx}
\end{table}

\paragraph{Note:}
IP-009 focuses on narrative AI systems. The SAP-OSS adaptation applies the stability monitoring principles (deviation measurement, discrete status indicators, governance metadata) to financial analytics while preserving the underlying governance-focused architecture.